%%%%%%%%%%%%%%%
% Remark 3: This section and subsections are mandatory - Please deliver a draft for assignments 1 and 2.
%%%%%%%%%%%%%%%

\section{Introduction}

\subsection{Motivation and Background}
As populations age, more older adults need care. Providing nursing-home care to everyone is expensive and often not what older adults prefer. Surveys and care-sector research from the United States, Switzerland, and the Netherlands indicate that most older adults prefer to remain in their own homes for as long as possible \cite{aarp2021aginginplace,unibas2020aginghome,rivm2026homecare}. To support this, new monitoring technologies are needed that allow older adults to remain independent while still enabling caregivers to intervene when assistance is needed. Continuous monitoring can provide information that occasional check-ins cannot, particularly for events requiring a rapid response. This may be especially useful for people with cognitive impairment or dementia, who can benefit from technologies that require little or no active interaction.

Existing monitoring technologies involve relevant trade-offs. Cameras provide detailed information about a person's location, posture, and movement, but also continuously record sensitive environments such as bedrooms, hospital rooms, and other private spaces. Wearable devices such as smartwatches avoid continuous visual monitoring, but require the user to wear, charge, and maintain the device. This dependence on user compliance can be particularly problematic for people with dementia, who may forget or refuse to wear such a device. Millimeter-wave (mmWave) radar offers a less intrusive alternative. Frequency-modulated continuous-wave (FMCW) radar measures reflected electromagnetic waves and can provide information about range, movement, and the spatial positions of objects without producing continuous video. Permanently installed radar sensors can therefore support passive monitoring without requiring active user interaction \cite{li2024sleepingpose,lai2023dual}.

For this project in particular, the motivation is to extend this type of general monitoring toward important aspects of everyday life and potentially dangerous situations. For us and Valtes, the important times of day are when a person is lying, sitting, or sleeping. Understanding whether the monitored person is in bed or on the toilet, whether they are asleep, which sleep stage they are in, and whether they are standing up or sitting down could provide useful signals for keeping the person happy, healthy, and safe.
